\section{AHSME 1981}
        \begin{enumerate}
        	\item (\textit{AHSME 1981}) If $\sqrt{x+2}=2$, then $(x+2)^{2}$ equals


 $\textbf{(A)}\ \sqrt{2}\qquad\textbf{(B)}\ 2\qquad\textbf{(C)}\ 4\qquad\textbf{(D)}\ 8\qquad\textbf{(E)}\ 16$


 

	\item (\textit{AHSME 1981}) Point $E$ is on side $AB$ of square $ABCD$. If $EB$ has length one and $EC$ has length two, then the area of the square is


 
                    \begin{figure}[H]
                        \centering
                        \includegraphics[width=0.5\textwidth]{Problem/AHSME_I_1981_2_Problem_Figure_1.png}
                    \end{figure}
                    


 

$\textbf{(A)}\ \sqrt{3}\qquad\textbf{(B)}\ \sqrt{5}\qquad\textbf{(C)}\ 3\qquad\textbf{(D)}\ 2\sqrt{3}\qquad\textbf{(E)}\ 5$


 

	\item (\textit{AHSME 1981}) For $x\neq0$, $\dfrac{1}{x}+\dfrac{1}{2x}+\dfrac{1}{3x}$ equals


 $\textbf{(A)}\ \dfrac{1}{2x}\qquad\textbf{(B)}\ \dfrac{1}{6x}\qquad\textbf{(C)}\ \dfrac{5}{6x}\qquad\textbf{(D)}\ \dfrac{11}{6x}\qquad\textbf{(E)}\ \dfrac{1}{6x^3}$


 

	\item (\textit{AHSME 1981}) If three times the larger of two numbers is four times the smaller and the difference between the numbers is $8$, the the larger of two numbers is 


 $\textbf{(A)}\ 16\qquad\textbf{(B)}\ 24\qquad\textbf{(C)}\ 32\qquad\textbf{(D)}\ 44\qquad\textbf{(E)}\ 52$


 

	\item (\textit{AHSME 1981}) In trapezoid $ABCD$, sides $AB$ and $CD$ are parallel, and diagonal $BD$ and side $AD$ have equal length. If $\angle DCB=110^\circ$ and $\angle CBD=30^\circ$, then $\angle ADB=$


 
                    \begin{figure}[H]
                        \centering
                        \includegraphics[width=0.5\textwidth]{Problem/AHSME_I_1981_5_Problem_Figure_1.png}
                    \end{figure}
                    


 $\textbf{(A)}\ 80^\circ\qquad\textbf{(B)}\ 90^\circ\qquad\textbf{(C)}\ 100^\circ\qquad\textbf{(D)}\ 110^\circ\qquad\textbf{(E)}\ 120^\circ$


 

	\item (\textit{AHSME 1981}) If $\dfrac{x}{x-1}=\dfrac{y^2+2y-1}{y^2+2y-2}$, then $x$ equals 


 $\textbf{(A)}\ y^2+2y-1\qquad\textbf{(B)}\ y^2+2y-2\qquad\textbf{(C)}\ y^2+2y+2 \qquad \\ \textbf{(D)}\ y^2+2y+1\qquad\textbf{(E)}\ -y^2-2y+1$


 

	\item (\textit{AHSME 1981}) How many of the first one hundred positive integers are divisible by all of the numbers $2$, $3$, $4$, and $5$?


 $\textbf{(A)}\ 0\qquad\textbf{(B)}\ 1\qquad\textbf{(C)}\ 2\qquad\textbf{(D)}\ 3\qquad\textbf{(E)}\ 4$


 

	\item (\textit{AHSME 1981}) For all positive numbers $x$, $y$, $z$, the product $(x+y+z)^{-1}(x^{-1}+y^{-1}+z^{-1})(xy+yz+xz)^{-1}[(xy)^{-1}+(yz)^{-1}+(xz)^{-1}]$ equals $\textbf{(A)}\ x^{-2}y^{-2}z^{-2}\qquad\textbf{(B)}\ x^{-2}+y^{-2}+z^{-2}\qquad\textbf{(C)}\ (x+y+z)^{-1}\qquad \textbf{(D)}\ \dfrac{1}{xyz}\qquad \\ \textbf{(E)}\ \dfrac{1}{xy+yz+xz}$


 

	\item (\textit{AHSME 1981}) In the adjoining figure, $PQ$ is a diagonal of the cube. If $PQ$ has length $a$, then the surface area of the cube is


 
                    \begin{figure}[H]
                        \centering
                        \includegraphics[width=0.5\textwidth]{Problem/AHSME_I_1981_9_Problem_Figure_1.png}
                    \end{figure}
                    


 $\textbf{(A)}\ 2a^2\qquad\textbf{(B)}\ 2\sqrt{2}a^2\qquad\textbf{(C)}\ 2\sqrt{3}a^2\qquad\textbf{(D)}\ 3\sqrt{3}a^2\qquad\textbf{(E)}\ 6a^2$


 

	\item (\textit{AHSME 1981}) The lines $L$ and $K$ are symmetric to each other with respect to the line $y=x$. If the equation of the line $L$ is $y=ax+b$ with $a\neq 0$ and $b\neq 0$, then the equation of $K$ is $y=$


 $\textbf{(A)}\ \dfrac{1}{a}x+b\qquad\textbf{(B)}\ -\dfrac{1}{a}x+b\qquad\textbf{(C)}\ \dfrac{1}{a}x-\dfrac{b}{a}\qquad\textbf{(D)}\ \dfrac{1}{a}x+\dfrac{b}{a}\qquad\textbf{(E)}\ \dfrac{1}{a}x-\dfrac{b}{a}$


 

	\item (\textit{AHSME 1981}) The three sides of a right triangle have integral lengths which form an arithmetic progression. One of the sides could have length


 $\textbf{(A)}\ 22\qquad\textbf{(B)}\ 58\qquad\textbf{(C)}\ 81\qquad\textbf{(D)}\ 91\qquad\textbf{(E)}\ 361$


 

	\item (\textit{AHSME 1981}) If $p$, $q$, and $M$ are positive numbers and $q<100$, then the number obtained by increasing $M$ by $p\%$ and decreasing the result by $q\%$ exceeds $M$ if and only if


 $\textbf{(A)}\ p>q \qquad\textbf{(B)}\ p>\dfrac{q}{100-q}\qquad\textbf{(C)}\ p>\dfrac{q}{1-q}\qquad \textbf{(D)}\ p>\dfrac{100q}{100+q}\qquad\textbf{(E)}\ p>\dfrac{100q}{100-q}$


 

	\item (\textit{AHSME 1981}) Suppose that at the end of any year, a unit of money has lost $10\%$ of the value it had at the beginning of that year. Find the smallest integer $n$ such that after $n$ years, the money will have lost at least $90\%$ of its value (To the nearest thousandth $\log_{10} 3 = 0.477$).


 $\textbf{(A)}\ 14\qquad\textbf{(B)}\ 16\qquad\textbf{(C)}\ 18\qquad\textbf{(D)}\ 20\qquad\textbf{(E)}\ 22$


 

	\item (\textit{AHSME 1981}) In a geometric sequence of real numbers, the sum of the first $2$ terms is $7$, and the sum of the first $6$ terms is $91$. The sum of the first $4$ terms is


 $\textbf{(A)}\ 28\qquad\textbf{(B)}\ 32\qquad\textbf{(C)}\ 35\qquad\textbf{(D)}\ 49\qquad\textbf{(E)}\ 84$


 

	\item (\textit{AHSME 1981}) If $b>1$, $x>0$, and $(2x)^{\log_b 2}-(3x)^{\log_b 3}=0$, then $x$ is


 $\textbf{(A)}\ \dfrac{1}{216}\qquad\textbf{(B)}\ \dfrac{1}{6}\qquad\textbf{(C)}\ 1\qquad\textbf{(D)}\ 6\qquad\textbf{(E)}\ \text{not uniquely determined}$


 

	\item (\textit{AHSME 1981}) The base three representation of $x$ is 

 \[12112211122211112222\]
The first digit (on the left) of the base nine representation of $x$ is


 $\textbf{(A)}\ 1\qquad\textbf{(B)}\ 2\qquad\textbf{(C)}\ 3\qquad\textbf{(D)}\ 4\qquad\textbf{(E)}\ 5$


 

	\item (\textit{AHSME 1981}) The function $f$ is not defined for $x=0$, but, for all non-zero real numbers $x$, $f(x)+2f\left(\dfrac{1}x\right)=3x$. The equation $f(x)=f(-x)$ is satisfied by


 $\textbf{(A)}\ \text{exactly one real number} \qquad \textbf{(B)}\ \text{exactly two real numbers} \qquad\textbf{(C)}\ \text{no real numbers}\qquad \\ \textbf{(D)}\ \text{infinitely many, but not all, non-zero real numbers} \qquad\textbf{(E)}\ \text{all non-zero real numbers}$


 

	\item (\textit{AHSME 1981}) The number of real solutions to the equation

 \[\dfrac{x}{100}=\sin x\]
is


 $\textbf{(A)}\ 61\qquad\textbf{(B)}\ 62\qquad\textbf{(C)}\ 63\qquad\textbf{(D)}\ 64\qquad\textbf{(E)}\ 65$


 

	\item (\textit{AHSME 1981}) In $\triangle ABC$, $M$ is the midpoint of side $BC$, $AN$ bisects $\angle BAC$, $BN\perp AN$, and $\theta$ is the measure of $\angle BAC$. If sides $AB$ and $AC$ have lengths $14$ and $19$, respectively, then find $MN$.


 
                    \begin{figure}[H]
                        \centering
                        \includegraphics[width=0.5\textwidth]{Problem/AHSME_I_1981_19_Problem_Figure_1.png}
                    \end{figure}
                    


 $\textbf{(A)}\ 2\qquad\textbf{(B)}\ \dfrac{5}{2}\qquad\textbf{(C)}\ \dfrac{5}{2}-\sin\theta\qquad\textbf{(D)}\ \dfrac{5}{2}-\dfrac{1}{2}\sin\theta\qquad\textbf{(E)}\ \dfrac{5}{2}-\dfrac{1}{2}\sin\left(\dfrac{1}{2}\theta\right)$


 

	\item (\textit{AHSME 1981}) A ray of light originates from point $A$ and travels in a plane, being reflected $n$ times between lines $AD$ and $CD$ before striking a point $B$ (which may be on $AD$ or $CD$) perpendicularly and retracing its path back to $A$ (At each point of reflection the light makes two equal angles as indicated in the adjoining figure. The figure shows the light path for $n=3$). If $\measuredangle CDA=8^\circ$, what is the largest value $n$ can have?


 
                    \begin{figure}[H]
                        \centering
                        \includegraphics[width=0.5\textwidth]{Problem/AHSME_I_1981_20_Problem_Figure_1.png}
                    \end{figure}
                    


 $\textbf{(A)}\ 6\qquad\textbf{(B)}\ 10\qquad\textbf{(C)}\ 38\qquad\textbf{(D)}\ 98\qquad\textbf{(E)}\ \text{There is no largest value.}$


 

	\item (\textit{AHSME 1981}) In a triangle with sides of lengths $a$, $b$, and $c$, $(a+b+c)(a+b-c) = 3ab$. The measure of the angle opposite the side length $c$ is


 $\textbf{(A)}\ 15^\circ\qquad\textbf{(B)}\ 30^\circ\qquad\textbf{(C)}\ 45^\circ\qquad\textbf{(D)}\ 60^\circ\qquad\textbf{(E)}\ 150^\circ$


 

	\item (\textit{AHSME 1981}) How many lines in a three dimensional rectangular coordinate system pass through four distinct points of the form $(i,j,k)$, where $i$, $j$, and $k$ are positive integers not exceeding four?


 $\textbf{(A)}\ 60\qquad\textbf{(B)}\ 64\qquad\textbf{(C)}\ 72\qquad\textbf{(D)}\ 76\qquad\textbf{(E)}\ 100$


 

	\item (\textit{AHSME 1981}) Equilateral $\triangle ABC$ is inscribed in a circle. A second circle is tangent internally to the circumcircle at $T$ and tangent to sides $AB$ and $AC$ at points $P$ and $Q$. If side $BC$ has length $12$, then segment $PQ$ has length


 
                    \begin{figure}[H]
                        \centering
                        \includegraphics[width=0.5\textwidth]{Problem/AHSME_I_1981_23_Problem_Figure_1.png}
                    \end{figure}
                    


 $\textbf{(A)}\ 6\qquad\textbf{(B)}\ 6\sqrt{3}\qquad\textbf{(C)}\ 8\qquad\textbf{(D)}\ 8\sqrt{3}\qquad\textbf{(E)}\ 9$


 

	\item (\textit{AHSME 1981}) If $\theta$ is a constant such that $0 < \theta < \pi$ and $x + \dfrac{1}{x} = 2\cos{\theta}$, then for each positive integer $n$, $x^n + \dfrac{1}{x^n}$ equals


 $\textbf{(A)}\ 2\cos\theta\qquad \textbf{(B)}\ 2^n\cos\theta\qquad \textbf{(C)}\ 2\cos^n\theta\qquad \textbf{(D)}\ 2\cos n\theta\qquad \textbf{(E)}\ 2^n\cos^n\theta$


 

	\item (\textit{AHSME 1981}) In $\triangle ABC$ in the adjoining figure, $AD$ and $AE$ trisect $\angle BAC$. The lengths of $BD$, $DE$ and $EC$ are $2$, $3$, and $6$, respectively. The length of the shortest side of $\triangle ABC$ is


 
                    \begin{figure}[H]
                        \centering
                        \includegraphics[width=0.5\textwidth]{Problem/AHSME_I_1981_25_Problem_Figure_1.png}
                    \end{figure}
                    


 $\textbf{(A)}\ 2\sqrt{10}\qquad \textbf{(B)}\ 11\qquad \textbf{(C)}\ 6\sqrt{6}\qquad \textbf{(D)}\ 6\qquad \textbf{(E)}\ \text{not uniquely determined by the given information}$


 

	\item (\textit{AHSME 1981}) Alice, Bob, and Carol repeatedly take turns tossing a die. Alice begins; Bob always follows Alice; Carol always follows Bob; and Alice always follows Carol. Find the probability that Carol will be the first one to toss a six. (The probability of obtaining a six on any toss is $\dfrac{1}{6}$, independent of the outcome of any other toss.)


 $\textbf{(A)}\ \dfrac{1}{3}\qquad \textbf{(B)}\ \dfrac{2}{9}\qquad \textbf{(C)}\ \dfrac{5}{18}\qquad \textbf{(D)}\ \dfrac{25}{91}\qquad \textbf{(E)}\ \dfrac{36}{91}$


 

	\item (\textit{AHSME 1981}) In the adjoining figure triangle $ABC$ is inscribed in a circle. Point $D$ lies on $\stackrel{\frown}{AC}$ with $\stackrel{\frown}{DC} = 30^\circ$, and point $G$ lies on $\stackrel{\frown}{BA}$ with $\stackrel{\frown}{BG}\, > \, \stackrel{\frown}{GA}$. Side $AB$ and side $AC$ each have length equal to the length of chord $DG$, and $\angle CAB = 30^\circ$. Chord $DG$ intersects sides $AC$ and $AB$ at $E$ and $F$, respectively. The ratio of the area of $\triangle AFE$ to the area of $\triangle ABC$ is


 
                    \begin{figure}[H]
                        \centering
                        \includegraphics[width=0.5\textwidth]{Problem/AHSME_I_1981_27_Problem_Figure_1.png}
                    \end{figure}
                    


 $\textbf{(A)}\ \dfrac {2 - \sqrt {3}}{3}\qquad \textbf{(B)}\ \dfrac {2\sqrt {3} - 3}{3}\qquad \textbf{(C)}\ 7\sqrt {3}-12\qquad \textbf{(D)}\ 3\sqrt {3}-5\qquad\\ \textbf{(E)}\ \dfrac {9-5\sqrt {3}}{3}$


 

	\item (\textit{AHSME 1981}) Consider the set of all equations $x^3 + a_2x^2 + a_1x + a_0 = 0$, where $a_2$, $a_1$, $a_0$ are real constants and $|a_i| \leq 2$ for $i = 0,1,2$. Let $r$ be the largest positive real number which satisfies at least one of these equations. Then


 $\textbf{(A)}\ 1 \leq r < \dfrac{3}{2}\qquad \textbf{(B)}\ \dfrac{3}{2} \leq r < 2\qquad \textbf{(C)}\ 2 \leq r < \dfrac{5}{2}\qquad \textbf{(D)}\ \dfrac{5}{2} \leq r < 3\qquad \textbf{(E)}\ 3 \leq r < \dfrac{7}{2}$


 

	\item (\textit{AHSME 1981}) If $a \geq 1$, then the sum of the real solutions of 


 $\sqrt{a - \sqrt{a + x}} = x$


 is equal to 


 $\textbf{(A)}\ \sqrt{a} - 1\qquad \textbf{(B)}\ \dfrac{\sqrt{a}- 1}{2}\qquad \textbf{(C)}\ \sqrt{a - 1}\qquad \textbf{(D)}\ \dfrac{\sqrt{a - 1}}{2}\qquad \textbf{(E)}\ \dfrac{\sqrt{4a- 3} - 1}{2}$


 

	\item (\textit{AHSME 1981}) If $a$, $b$, $c$, and $d$ are the solutions of the equation $x^4 - bx - 3 = 0$, then an equation whose solutions are 

 \[\dfrac {a + b + c}{d^2}, \dfrac {a + b + d}{c^2}, \dfrac {a + c + d}{b^2}, \dfrac {b + c + d}{a^2}\]is


 $\textbf{(A)}\ 3x^4 + bx + 1 = 0\qquad \textbf{(B)}\ 3x^4 - bx + 1 = 0\qquad \textbf{(C)}\ 3x^4 + bx^3 - 1 = 0\qquad \\\textbf{(D)}\ 3x^4 - bx^3 - 1 = 0\qquad  \textbf{(E)}\ \text{none of these}$


 

\end{enumerate}